% J09 — A Small Commutative Non-Associative Magma on Z/10Z with
%        Role-Deterministic Boundary Behavior
%
% Brayden R. Sanders, M. Gish (2026)
% Submitted to: Algebra Universalis
% MSC 2020: 20N02, 08A05, 20N05
% License (this paper + bundled tables): CC-BY-4.0

\documentclass[11pt,reqno]{amsart}

\usepackage{amsmath, amssymb, amsthm, mathtools}
\usepackage[margin=1in]{geometry}
\usepackage[colorlinks=true,linkcolor=blue,citecolor=blue,urlcolor=blue]{hyperref}
\usepackage{microtype}
\usepackage{booktabs}
\usepackage{array}

\theoremstyle{plain}
\newtheorem{theorem}{Theorem}[section]
\newtheorem{proposition}[theorem]{Proposition}
\newtheorem{lemma}[theorem]{Lemma}
\newtheorem{corollary}[theorem]{Corollary}

\theoremstyle{definition}
\newtheorem{definition}[theorem]{Definition}

\theoremstyle{remark}
\newtheorem*{remark}{Remark}

\newcommand{\Z}{\mathbb{Z}}
\newcommand{\F}{\mathbb{F}}
\newcommand{\TS}{\mathrm{TSML}}
\newcommand{\BH}{\mathrm{BHML}}
\newcommand{\sig}{\sigma}
\newcommand{\Zten}{\Z/10\Z}

\title[A small role-deterministic magma on $\Zten$]
{A Small Commutative Non-Associative Magma on $\Zten$\\
 with Role-Deterministic Boundary Behavior}

\author{Brayden R.\ Sanders}
\address{7Site LLC, Hot Springs, Arkansas, USA}
\email{brayden@7site.co}

\author{M.\ Gish}
\address{Independent Researcher, Hot Springs, Arkansas, USA}
\email{monica.gish1992@gmail.com}

\subjclass[2020]{20N02, 08A05, 20N05}
\keywords{commutative magma, finite non-associative algebra,
$\Z/10\Z$, role partition, semi-factorization, role-deterministic
boundary, Dr\'apal--Wanless neighborhood}

\date{\today}

\begin{document}

\begin{abstract}
We study a small commutative non-associative magma $M_R$ on the
four-element set $\{V, F, S, T\}$, obtained as the role-mode
reduction of a specific commutative binary operation $\BH$ on
$\Zten$ along the four-block partition
$F = \{1, 3, 5, 7, 9\}$, $S = \{2, 4, 8\}$, $T = \{6\}$,
$V = \{0\}$. The operation $\BH$ is given explicitly by an
elementary four-zone rule (Definition~\ref{def:BH}), so that the
paper is self-contained: the reader can verify every claim by
inspecting the $10 \times 10$ table.

We prove: (i) $M_R$ is commutative, has $V$ as two-sided identity,
and is non-associative, with explicit witness
$M_R(M_R(F, F), S) = F \neq T = M_R(F, M_R(F, S))$
(Theorems~\ref{thm:V-identity} and \ref{thm:role-magma-non-assoc}).
(ii) $\BH$ is \emph{role-deterministic on the boundary}: for every
input pair $(a, b) \in \Zten \times \Zten$ with at least one of
$a, b$ in $V \cup T$, the output role is determined by the input
roles alone, independently of the operator-level identity of $a$
and $b$. This determinism fails on the role-pairs in $\{F, S\}^2$,
where the output role distribution is non-trivial; we list the
exact distributions in Theorem~\ref{thm:semi-fact}. (iii) The
$\BH$-self-orbit period function on $\{1, \ldots, 6\}$ is the
linear function $\mathrm{period}(n) = 7 - n$, with boundary
values $4, 3, 2, 1$ at $n = 7, 8, 9, 0$
(Lemma~\ref{lem:bh-diagonal-period}).

As a side observation, we record one $\Z$-valued row statistic
$\Psi(n)$ on the pair of tables $(\TS, \BH)$ whose row sum is
$\Psi(\Zten) = 21 = T_6$, the sixth triangular number, and we
record its decomposition along $\sig$-fixed and $\sig$-cycle
index sets ($\{-1, 22\}$). The integer $21$ is not interpreted
here; we merely verify the row sum and note the orbit-type
decomposition.

We do not claim that $M_R$ exhibits a knot-theoretic structure
or that it implements a Crossing Lemma or a Rademacher invariant.
Such suggestive scaffolding appears in adjacent companion work
\cite{FourCore} and is not invoked here. The present paper is a
study of $M_R$ as a small finite commutative non-associative
algebra in the Dr\'apal--Wanless~\cite{DrapalWanless2021}
neighborhood, with role-deterministic boundary behavior as its
distinguishing property.
\end{abstract}

\maketitle

%==============================================================
\section*{0. Lens, substrate, and tier discipline}
%==============================================================

\noindent
\textsc{Lens and substrate.}
The substrate is $\Zten$ with two specific commutative binary
tables $\BH$ and $\TS$ defined explicitly by elementary rules in
Section~\ref{sec:tables}; the four-block role partition
$F, S, T, V$ is a labeling-by-fiat motivated by an interpretive
reading not used in the proofs. Whether some other small
commutative magma on $\Zten$ admits an analogous role-mode
reduction is open; the present paper records the structure of
$M_R$ for the specific choice $\BH$.

\medskip\noindent
\textsc{Tier discipline.}
\begin{itemize}\setlength\itemsep{2pt}
\item \emph{PROVEN.} The role-magma $M_R$ is commutative, has
$V$ as two-sided identity, and is non-associative
(Theorems~\ref{thm:V-identity}, \ref{thm:role-magma-non-assoc}).
$\BH$ is role-deterministic on every input pair containing $V$
or $T$, with the exact role-output distributions on
$\{F, S\}^2$ given in Theorem~\ref{thm:semi-fact}. The
$\BH$-self-orbit period function on $\{1, \ldots, 6\}$ is
linear, $\mathrm{period}(n) = 7 - n$
(Lemma~\ref{lem:bh-diagonal-period}).
\item \emph{COMPUTED.} The script
\texttt{verify\_role\_magma.py} (bundled with the manuscript)
verifies the $4 \times 4$ role-mode table from $\BH$ in
$<0.1$\,s; it also re-derives the period function
on $\{0, 1, \ldots, 9\}$ and the row-asymmetry computation
of Section~\ref{sec:row-asymmetry}.
\item \emph{STRUCTURAL RHYME.} The row-asymmetry function $\Psi$
of Section~\ref{sec:row-asymmetry} sums to $\Psi(\Zten) = 21
= T_6$ (the sixth triangular number) on the full $\Zten$, and
splits as $\Psi(\sig\text{-fixed}) + \Psi(\sig\text{-cycle}) =
-1 + 22 = 21$ along $\sig$-orbit type. We record the
total and the orbit decomposition as numerical observations;
neither is asserted as a structural theorem.
\item \emph{OPEN.} A characterization of the role partition
$\{V, F, S, T\}$ as a congruence of some derived structure
(it is not a congruence of $\BH$ itself). Whether
analogous role-deterministic-boundary algebras exist on
$\Z/14\Z$ or $\Z/18\Z$ is open. The closest published
precedent is the
Dr\'apal--Wanless~\cite{DrapalWanless2021} neighborhood
of small finite commutative non-associative structures.
\end{itemize}

%==============================================================
\section{The two binary operations on $\Zten$}
\label{sec:tables}
%==============================================================

We define two commutative binary operations on
$\Zten = \{0, 1, \ldots, 9\}$ explicitly. The reader can verify
every claim of the paper by inspecting the $10 \times 10$ tables
in Tables~\ref{tab:BH} and \ref{tab:TS} below.

\begin{definition}[The operation $\BH$]\label{def:BH}
The commutative binary operation $\BH : \Zten \times \Zten
\to \Zten$ is defined by the four-zone rule:
\begin{enumerate}\setlength\itemsep{2pt}
\item[\rm (R$_A$)] \textbf{$0$-row and $0$-column.} For all $j$,
$\BH(0, j) = j$ and $\BH(i, 0) = i$.
\item[\rm (R$_B$)] \textbf{Core zone.} For
$i, j \in \{1, \ldots, 6\}$, $\BH(i, j) = \max(i, j) + 1$.
\item[\rm (R$_7$)] \textbf{$7$-row and $7$-column.} For
$j \ge 1$, $\BH(7, j) = (j + 1) \bmod 10$, and symmetrically
$\BH(i, 7) = (i + 1) \bmod 10$ for $i \ge 1$.
\item[\rm (R$_{89}$)] \textbf{Boundary zone.} The remaining
$32$ cells (with $i \in \{8, 9\}$ or $j \in \{8, 9\}$ and
the prior zones excluded) take the values listed
explicitly in Table~\ref{tab:BH}.
\end{enumerate}
\end{definition}

\begin{table}[h]
\caption{The full $\BH$ table, $\BH(i, j) =$ entry at row $i$,
column $j$. Rows and columns indexed $0, 1, \ldots, 9$.
Symmetric: $\BH(i, j) = \BH(j, i)$.}
\label{tab:BH}
\centering
\begin{tabular}{c|cccccccccc}
$\BH$ & $0$ & $1$ & $2$ & $3$ & $4$ & $5$ & $6$ & $7$ & $8$ & $9$ \\ \hline
$0$ & $0$ & $1$ & $2$ & $3$ & $4$ & $5$ & $6$ & $7$ & $8$ & $9$ \\
$1$ & $1$ & $2$ & $3$ & $4$ & $5$ & $6$ & $7$ & $2$ & $6$ & $6$ \\
$2$ & $2$ & $3$ & $3$ & $4$ & $5$ & $6$ & $7$ & $3$ & $6$ & $6$ \\
$3$ & $3$ & $4$ & $4$ & $4$ & $5$ & $6$ & $7$ & $4$ & $6$ & $6$ \\
$4$ & $4$ & $5$ & $5$ & $5$ & $5$ & $6$ & $7$ & $5$ & $7$ & $7$ \\
$5$ & $5$ & $6$ & $6$ & $6$ & $6$ & $6$ & $7$ & $6$ & $7$ & $7$ \\
$6$ & $6$ & $7$ & $7$ & $7$ & $7$ & $7$ & $7$ & $7$ & $7$ & $7$ \\
$7$ & $7$ & $2$ & $3$ & $4$ & $5$ & $6$ & $7$ & $8$ & $9$ & $0$ \\
$8$ & $8$ & $6$ & $6$ & $6$ & $7$ & $7$ & $7$ & $9$ & $7$ & $8$ \\
$9$ & $9$ & $6$ & $6$ & $6$ & $7$ & $7$ & $7$ & $0$ & $8$ & $0$ \\
\end{tabular}
\end{table}

\begin{definition}[The operation $\TS$]\label{def:TS}
The commutative binary operation $\TS : \Zten \times \Zten
\to \Zten$ is defined by:
\begin{enumerate}\setlength\itemsep{2pt}
\item[\rm (V0)] \textbf{Row $0$.} For all $j$, $\TS(0, j) = 0$
if $j \neq 7$, and $\TS(0, 7) = 7$.
\item[\rm (V1)] \textbf{Column $0$.} Symmetric to (V0).
\item[\rm (E)] \textbf{Echo cells.} The five symmetric pairs
$\{(1,2), (2,4), (2,9), (3,9), (4,8)\}$ together with their
transposes take the values listed in Table~\ref{tab:TS}.
\item[\rm (D)] \textbf{Default.} Every cell not in (V0), (V1),
(E) takes the value $7$.
\end{enumerate}
\end{definition}

\begin{table}[h]
\caption{The full $\TS$ table. Symmetric: $\TS(i, j) = \TS(j, i)$.}
\label{tab:TS}
\centering
\begin{tabular}{c|cccccccccc}
$\TS$ & $0$ & $1$ & $2$ & $3$ & $4$ & $5$ & $6$ & $7$ & $8$ & $9$ \\ \hline
$0$ & $0$ & $0$ & $0$ & $0$ & $0$ & $0$ & $0$ & $7$ & $0$ & $0$ \\
$1$ & $0$ & $7$ & $3$ & $7$ & $7$ & $7$ & $7$ & $7$ & $7$ & $7$ \\
$2$ & $0$ & $3$ & $7$ & $7$ & $4$ & $7$ & $7$ & $7$ & $7$ & $9$ \\
$3$ & $0$ & $7$ & $7$ & $7$ & $7$ & $7$ & $7$ & $7$ & $7$ & $3$ \\
$4$ & $0$ & $7$ & $4$ & $7$ & $7$ & $7$ & $7$ & $7$ & $8$ & $7$ \\
$5$ & $0$ & $7$ & $7$ & $7$ & $7$ & $7$ & $7$ & $7$ & $7$ & $7$ \\
$6$ & $0$ & $7$ & $7$ & $7$ & $7$ & $7$ & $7$ & $7$ & $7$ & $7$ \\
$7$ & $7$ & $7$ & $7$ & $7$ & $7$ & $7$ & $7$ & $7$ & $7$ & $7$ \\
$8$ & $0$ & $7$ & $7$ & $7$ & $8$ & $7$ & $7$ & $7$ & $7$ & $7$ \\
$9$ & $0$ & $7$ & $9$ & $3$ & $7$ & $7$ & $7$ & $7$ & $7$ & $7$ \\
\end{tabular}
\end{table}

\subsection*{The permutation $\sig$.}

We define the permutation
\[
  \sig = (0)(3)(8)(9)(1\,7\,6\,5\,4\,2) \in S_{10},
\]
a single $6$-cycle on $\{1, 2, 4, 5, 6, 7\}$ together with four
fixed points on $\{0, 3, 8, 9\}$. Note that $\sig$ has order
$6$, not $2$; hence $\sig$ is a permutation, not an
involution. We use the explicit involution
\[
  \sig^3 = (1\,5)(2\,6)(4\,7)(0)(3)(8)(9)
\]
(the cube of $\sig$, of order $2$) at one place in
Section~\ref{sec:row-asymmetry}; otherwise we use the
$\sig$-orbit structure directly.

\subsection*{Closest published precedent.}

The two operations $\BH$ and $\TS$ are commutative non-associative
magmas on a $10$-element set with explicit elementary rules. The
closest published precedent for this neighborhood is
Dr\'apal--Wanless~\cite{DrapalWanless2021}, on maximally
non-associative quasigroups of small order. Both $\BH$ and $\TS$
are far from maximally non-associative; both are commutative
(Dr\'apal--Wanless 2021's quasigroups need not be); both
preserve the $4$-element subset $\{0, 7, 8, 9\}$ (a load-bearing
property of companion work~\cite{FourCore}). The intellectual
neighborhood is the same; the specific structural conditions
differ.

%==============================================================
\section{The role partition and the role magma $M_R$}
\label{sec:role-magma}
%==============================================================

\subsection*{The four-block partition.}

We partition $\Zten = \{0, 1, \ldots, 9\}$ into four blocks:
\[
  V := \{0\}, \quad F := \{1, 3, 5, 7, 9\}, \quad
  S := \{2, 4, 8\}, \quad T := \{6\}.
\]
Sizes $|V| = 1$, $|F| = 5$, $|S| = 3$, $|T| = 1$. The partition
is a labeling-by-fiat motivated by interpretive reading; we do
\emph{not} claim it is a congruence of $\BH$ (it is not, because
$\BH$ is not coarsenable to $\{V, F, S, T\}$ uniformly; see
Theorem~\ref{thm:semi-fact} below for the failure on
$\{F, S\}^2$).

We write $\rho : \Zten \to \{V, F, S, T\}$ for the role function:
$\rho(0) = V$, $\rho(1) = \rho(3) = \rho(5) = \rho(7) = \rho(9)
= F$, $\rho(2) = \rho(4) = \rho(8) = S$, $\rho(6) = T$.

\subsection*{The role-mode reduction.}

For each input role-pair $(r_a, r_b) \in \{V, F, S, T\}^2$, we
form the multiset of output roles
\[
  \mathcal{O}(r_a, r_b) := \{\rho(\BH(a, b)) : a \in \rho^{-1}(r_a),
  \ b \in \rho^{-1}(r_b)\}
\]
and define $M_R(r_a, r_b)$ as the \emph{mode} (most frequent
element) of this multiset. The full role-mode table is:

\begin{equation*}
\begin{array}{c|cccc}
M_R & V & F & S & T \\ \hline
V & V & F & S & T \\
F & F & T & F & F \\
S & S & F & F & F \\
T & T & F & F & F
\end{array}
\end{equation*}

The output-role distributions on $\{F, S\}^2$ are:
\begin{itemize}\setlength\itemsep{2pt}
\item $\mathcal{O}(F, F) = \{F: 2,\ S: 9,\ T: 11,\ V: 3\}$ (size
$25$, mode $T$).
\item $\mathcal{O}(F, S) = \mathcal{O}(S, F) = \{T: 5,\ F: 8,\
S: 2\}$ (size $15$, mode $F$).
\item $\mathcal{O}(S, S) = \{F: 7,\ T: 2\}$ (size $9$, mode $F$).
\end{itemize}

The remaining nine role-pairs (containing $V$ or $T$) have
\emph{singleton} output distributions: $\mathcal{O}(V, V) =
\{V\}$, $\mathcal{O}(V, F) = \{F: 5\}$, $\mathcal{O}(V, S) =
\{S: 3\}$, $\mathcal{O}(V, T) = \{T\}$, $\mathcal{O}(F, T) =
\{F: 5\}$, $\mathcal{O}(S, T) = \{F: 3\}$, $\mathcal{O}(T, T)
= \{F\}$, plus the symmetric mirrors. In each of these nine
cases, the output role is uniquely determined by the role-pair
of inputs.

\subsection{Properties of $M_R$}

\begin{theorem}[$V$ is a two-sided identity in $M_R$]
\label{thm:V-identity}
For every $r \in \{V, F, S, T\}$,
$M_R(V, r) = M_R(r, V) = r$. Moreover,
$M_R(V, V) = V$ is the unique idempotent.
\end{theorem}

\begin{proof}
By direct inspection of the role-mode table: $M_R(V, V) = V$,
$M_R(V, F) = F$, $M_R(V, S) = S$, $M_R(V, T) = T$. Symmetric for
the column. The unique idempotent: from the $4 \times 4$ table,
the cells with $M_R(r, r) = r$ are $r = V$ ($M_R(V, V) = V$);
$r = F$ has $M_R(F, F) = T \neq F$; $r = S$ has $M_R(S, S) = F
\neq S$; $r = T$ has $M_R(T, T) = F \neq T$. So $V$ is the unique
idempotent.
\end{proof}

\begin{theorem}[Commutativity, non-associativity]
\label{thm:role-magma-non-assoc}
$M_R$ is commutative: $M_R(r_a, r_b) = M_R(r_b, r_a)$ for all
$r_a, r_b$. $M_R$ is non-associative; an explicit witness is
\[
  M_R(M_R(F, F), S)
  = M_R(T, S) = F
  \neq T = M_R(F, F) = M_R(F, M_R(F, S)).
\]
\end{theorem}

\begin{proof}
Commutativity: by direct inspection of the role-mode table, the
matrix is symmetric. (Equivalently: $\BH$ is commutative, so the
output-role multiset $\mathcal{O}(r_a, r_b)$ is invariant under
swapping the inputs.)

Non-associativity: from the role-mode table,
$M_R(F, F) = T$, $M_R(T, S) = F$. So
$M_R(M_R(F, F), S) = M_R(T, S) = F$. Separately,
$M_R(F, S) = F$, so $M_R(F, M_R(F, S)) = M_R(F, F) = T$. Since
$F \neq T$, $M_R$ is non-associative.
\end{proof}

\subsection{Role-deterministic boundary, branching interior}

\begin{theorem}[Role-deterministic boundary]
\label{thm:semi-fact}
Let $a, b \in \Zten$ with $\rho(a), \rho(b) \in \{V, F, S, T\}$.
The output role $\rho(\BH(a, b))$ depends only on the input
roles $\rho(a), \rho(b)$ when at least one of $\rho(a), \rho(b)$
is $V$ or $T$. Concretely:
\begin{enumerate}\setlength\itemsep{2pt}
\item For every $(r_a, r_b) \in \{V, F, S, T\}^2$ with
$V \in \{r_a, r_b\}$ or $T \in \{r_a, r_b\}$, the multiset
$\mathcal{O}(r_a, r_b)$ is supported on a single role
(``role-deterministic''):
\begin{align*}
\mathcal{O}(V, V) &= \{V\}, \quad
\mathcal{O}(V, F) = \{F : 5\}, \quad
\mathcal{O}(V, S) = \{S : 3\}, \\
\mathcal{O}(V, T) &= \{T\}, \quad
\mathcal{O}(F, T) = \{F : 5\}, \quad
\mathcal{O}(S, T) = \{F : 3\}, \quad
\mathcal{O}(T, T) = \{F\},
\end{align*}
plus the symmetric pairs.

\item For $(r_a, r_b) \in \{F, S\}^2$, the multiset
$\mathcal{O}(r_a, r_b)$ is non-trivially supported on at least
two roles (``role-branching''), with the explicit distributions
$\mathcal{O}(F, F)$, $\mathcal{O}(F, S) = \mathcal{O}(S, F)$,
$\mathcal{O}(S, S)$ given immediately above.
\end{enumerate}
\end{theorem}

\begin{proof}
Direct verification on the table $\BH$ in Table~\ref{tab:BH},
restricted to the rows/columns prescribed by each role-pair.
The script \texttt{verify\_role\_magma.py} bundled with the
manuscript carries out the enumeration in $<0.1$\,s.

For the $V$-row and $V$-column the $\BH$ rules (R$_A$) give
$\BH(0, j) = j$ and $\BH(i, 0) = i$, so the output role equals
the role of the non-$V$ input; this gives the four $V$-distributions.
For $T = \{6\}$, the row $\BH(6, j)$ is $\{6, 7, 7, 7, 7, 7, 7,
7, 7, 7\}$, so $\BH(T, F)$ takes value $7 \in F$ for the five
$F$-inputs $j \in \{1, 3, 5, 7, 9\}$ (giving $\mathcal{O}(T, F)
= \{F : 5\}$); for $S$-inputs $j \in \{2, 4, 8\}$, the values are
$\{7, 7, 7\}$, giving $\mathcal{O}(T, S) = \{F : 3\}$; for the
single $T$-input $j = 6$, $\BH(6, 6) = 7$, giving
$\mathcal{O}(T, T) = \{F\}$. The $V \cdot T$ cell is $\BH(0, 6)
= 6 \in T$, giving $\mathcal{O}(V, T) = \{T\}$.

For the $\{F, S\}^2$ block: by symmetric enumeration over the
$5 \cdot 5 = 25$ cells in $F \times F$, $5 \cdot 3 = 15$ cells in
$F \times S$, and $3 \cdot 3 = 9$ cells in $S \times S$, we get
the listed role-output distributions. These are non-trivial
multi-element multisets, hence the failure of role-determinism on
$\{F, S\}^2$.
\end{proof}

\begin{remark}[Reading: a $4$-block partition of $\Zten$ that is
not a congruence of $\BH$]
The role partition $\{V, F, S, T\}$ is a congruence of $\BH$
restricted to inputs containing $V$ or $T$, but it fails to be a
congruence on $\{F, S\}^2$. The role-mode reduction $M_R$
``forgets'' the role-multiplicities on $\{F, S\}^2$, projecting
each multiset to its mode. This is a lossy reduction: the
distribution $\{F : 2, S : 9, T : 11, V : 3\}$ is not recoverable
from the mode value $T$ alone. The role-mode reduction $M_R$ is
therefore not a quotient algebra in the sense of universal
algebra; it is a small algebra obtained from $\BH$ by a
mode-projection, and it inherits commutativity and an identity
element, but not associativity.
\end{remark}

%==============================================================
\section{The $\BH$-self-orbit period function}
\label{sec:periods}
%==============================================================

For each $n \in \Zten$, define the $\BH$-self-orbit by iterated
diagonal: $n_0 := n$, $n_{k+1} := \BH(n_k, n_k)$. Define the
\emph{period} of $n$ as the smallest $p \ge 1$ such that the
sequence $(n_k)_{k \ge 0}$ enters a cycle of length $p$, with
$\mathrm{period}(n) := p$.

\begin{lemma}[Linear first-passage time to $7$ on
$\{1, \ldots, 6\}$]
\label{lem:bh-diagonal-period}
Let $\tau(n) := \min\{k \ge 0 : n_k = 7\}$ denote the
first-passage time of the $\BH$-self-orbit at $n$ to the value
$7$, with the convention $\tau(n) = \infty$ if no such $k$
exists. Then
\[
  \tau(n) = 7 - n \qquad \text{for } n \in \{1, 2, \ldots, 7\},
\]
$\tau(8) = 1$, and $\tau(0) = \tau(9) = \infty$ (the orbits at
$0$ and $9$ both fall into the absorbing fixed point $0$ and
never reach $7$).
\end{lemma}

\begin{proof}
By direct iteration on the diagonal of $\BH$ (Table~\ref{tab:BH}).
The diagonal values are $\BH(0, 0) = 0$, $\BH(n, n) = n + 1$ for
$n \in \{1, \ldots, 6\}$, $\BH(7, 7) = 8$, $\BH(8, 8) = 7$,
$\BH(9, 9) = 0$. The orbit at $n \in \{1, \ldots, 6\}$ is
$n, n + 1, n + 2, \ldots, 6, 7$, hitting $7$ at step $7 - n$.
The orbit at $n = 7$ has $\tau(7) = 0$ trivially. The orbit at
$n = 8$: $\BH(8, 8) = 7$, so $\tau(8) = 1$. The orbit at
$n = 9$: $\BH(9, 9) = 0$, then $\BH(0, 0) = 0$, so the
orbit becomes $9, 0, 0, \ldots$ and never reaches $7$, giving
$\tau(9) = \infty$. Similarly $\tau(0) = \infty$.
\end{proof}

\begin{remark}[Cycle structure on $\{1, \ldots, 8\}$]
\label{rem:cycle-78}
The eventual cycle structure of the $\BH$-self-orbit dynamics is:
each input $n \in \{1, 2, \ldots, 8\}$ eventually enters the
$2$-cycle $\{7 \leftrightarrow 8\}$; the inputs $\{0, 9\}$
instead fall into the fixed-point $\{0\}$. The dynamics
therefore have two attractor basins (a $2$-cycle and a
fixed-point), and the linear first-passage time
$\tau(n) = 7 - n$ on $\{1, \ldots, 6\}$ is a natural statistic
of the dominant basin's transient.
\end{remark}

%==============================================================
\section{Row-asymmetry and the integer $21$}
\label{sec:row-asymmetry}
%==============================================================

We compute one $\Z$-valued row-statistic of the pair $(\TS, \BH)$
whose row sum is $21 = T_6$, the sixth triangular number.

\begin{definition}[Row-asymmetry function]\label{def:row-asym}
For $n \in \Zten$, let
\[
  \Psi(n) := \#\{j \in \Zten : \TS(n, j) > \BH(n, j)\}
            - \#\{j \in \Zten : \BH(n, j) > \TS(n, j)\}.
\]
\end{definition}

By direct computation from Tables~\ref{tab:BH} and \ref{tab:TS}:
\begin{center}
\begin{tabular}{c|cccccccccc|c}
$n$ & $0$ & $1$ & $2$ & $3$ & $4$ & $5$ & $6$ & $7$ & $8$ & $9$ & total \\ \hline
$\Psi(n)$ & $-8$ & $6$ & $4$ & $5$ & $4$ & $5$ & $-1$ & $4$ & $1$ & $1$ & $21$ \\
\end{tabular}
\end{center}

(The row-by-row counts are reproduced by
\texttt{verify\_role\_magma.py}; the total $\sum_n \Psi(n) = 21$
is a finite check.)

\begin{proposition}[$\sig$-orbit decomposition]
\label{prop:triangular}
Let $\Psi$ be the row-asymmetry function of
Definition~\ref{def:row-asym}. Then
\[
  \sum_{n \in \{0,3,8,9\}} \Psi(n)
  = -8 + 5 + 1 + 1 = -1,
  \qquad
  \sum_{n \in \{1, 2, 4, 5, 6, 7\}} \Psi(n)
  = 6 + 4 + 4 + 5 - 1 + 4 = 22,
\]
so $\Psi(\text{$\sig$-fixed indices}) + \Psi(\text{$\sig$-cycle
indices}) = -1 + 22 = 21 = T_6$, the sixth triangular number.
\end{proposition}

\begin{proof}
Direct addition of the explicit $\Psi$-values listed above.
\end{proof}

\begin{remark}[On stronger numerical patterns]\label{rem:no-fibonacci}
We considered finer decompositions of the integer $21$ along the
$\sig$-orbit structure or the role partition (e.g., a Fibonacci
or fine-triangular decomposition). The $\Psi$-values do
\emph{not} reproduce a clean Fibonacci or fine-triangular
decomposition on the full $\Zten$:
\begin{align*}
\sum_{n \in F} \Psi(n) &= \Psi(1) + \Psi(3) + \Psi(5) + \Psi(7)
+ \Psi(9) = 6 + 5 + 5 + 4 + 1 = 21, \\
\sum_{n \in S} \Psi(n) &= \Psi(2) + \Psi(4) + \Psi(8) = 4 + 4 + 1
= 9, \\
\sum_{n \in T} \Psi(n) &= \Psi(6) = -1, \\
\sum_{n \in V} \Psi(n) &= \Psi(0) = -8.
\end{align*}
These do not match the Fibonacci pair $(F_7, F_6) = (13, 8)$ or
any structurally clean sum-decomposition we have identified. We
record the row-sum total $\Psi(\Zten) = 21 = T_6$ and the
$\sig$-orbit decomposition $\{-1, 22\}$ as the only structurally
robust observations of this section. Whether $\Psi$ admits a
finer Fibonacci-type decomposition under a sign-and-convention
adjustment we have not explored is an open question, related to
the period-to-trace lifts discussed in adjacent companion
notes~\cite{FourCore}; it is not invoked or used in the present
paper.
\end{remark}

%==============================================================
\section{Verification}
%==============================================================

A short Python script
\texttt{verify\_role\_magma.py} (bundled with the manuscript)
verifies, in under $0.1$ second on a standard laptop:
\begin{itemize}\setlength\itemsep{2pt}
\item The $4 \times 4$ role-mode table $M_R$ from $\BH$ via the
explicit role partition (Theorem~\ref{thm:V-identity},
\ref{thm:role-magma-non-assoc}).
\item The role-output multiset distributions
$\mathcal{O}(r_a, r_b)$ for all $16$ role-pairs
(Theorem~\ref{thm:semi-fact}).
\item The $\BH$-self-orbit period function on
$\{0, 1, \ldots, 9\}$ (Lemma~\ref{lem:bh-diagonal-period}).
\item The row-asymmetry function $\Psi$ and its
$\sig$-orbit-type and role-partition decompositions
(Section~\ref{sec:row-asymmetry}).
\end{itemize}

The script imports the tables $\BH$ and $\TS$ from a self-contained
local module \texttt{ck\_tables.py}, also bundled, licensed under
CC-BY-4.0. No external data dependencies.

%==============================================================
\section{What this paper does not claim}
%==============================================================

We make explicit the limits of the present paper:
\begin{itemize}\setlength\itemsep{2pt}
\item We do \emph{not} claim that $M_R$ exhibits a knot-theoretic
structure. The companion working paper~\cite{FourCore} discusses
analogies with the modular surface
\cite{KatokUgarcovici2007}, with knots-and-primes
\cite{Morishita2024}, and with Rademacher symbols
\cite{MatsusakaUeki2023, Burrin2024}; these analogies are
\emph{conceptual scaffolding}, not theorems of the present paper.
\item We do \emph{not} claim that the integer $21$ has a
PSL$(2, \Z)$-Rademacher interpretation. Five lifts to PSL$(2,
\Z)$ words from the $\BH$-self-orbits have been tested and none
produces $\pm 21$ as an exact group-theoretic invariant. The
period-to-trace bridge under a chosen $2 \times 2$ representative
gives $-21$ numerically, but it is one choice of lift among many
and we do not derive it.
\item We do \emph{not} introduce a class of ``paradoxical
information algebras'' as defined algebraic objects. The phrase
appears in some companion notes as descriptive language for the
role-deterministic-vs-branching dichotomy of
Theorem~\ref{thm:semi-fact}; defining a class of algebraic
structures sharing this property is left to future work and is
\emph{not} done here.
\item We do \emph{not} introduce a ``trefoil characterization''
predicated on a runtime processor. Earlier internal drafts
contained such a characterization, but the predicate
``trefoil-equivalent'' depends on a computational assignment we
do not make formal in the present paper. The relevant section
has been excised.
\end{itemize}

%==============================================================
\section{Open questions}
%==============================================================

\begin{enumerate}\setlength\itemsep{2pt}
\item Is the role partition $\{V, F, S, T\}$ a congruence of any
derived structure (e.g., a $\BH$-iterated derivative, a
restriction to a sub-magma)? The role-mode reduction $M_R$ is a
finite commutative magma; whether it embeds canonically into
$\BH$ via a quotient construction is open.
\item Does an analogous role-mode reduction exist for $\BH$
restricted to a 4-element subset like $\{0, 7, 8, 9\}$?
\item Do analogous role-deterministic-boundary algebras exist on
$\Z/14\Z$, $\Z/12\Z$, $\Z/8\Z$? The companion paper~\cite{FourCore}
records ring-extension scaling for the joint
$\{V, H, B, R\}$-preservation property; whether the
role-deterministic-boundary property scales similarly is not
addressed here.
\end{enumerate}

%==============================================================
\section*{Acknowledgments}
%==============================================================

This work is a self-contained algebraic-combinatorial study;
familiarity with adjacent companion work~\cite{FourCore} is not
required.

%==============================================================
\begin{thebibliography}{99}
%==============================================================

\bibitem{DrapalWanless2021}
A.~Dr\'apal, I.~M.~Wanless, \emph{Maximally non-associative
quasigroups}, J.\ Combin.\ Theory Ser.\ A \textbf{184} (2021),
105510.

\bibitem{FourCore} B.~R.~Sanders, M.~Gish,
\emph{Joint Closure, Per-Coordinate Fuse Data, and a Closed-Form
Algebraic Attractor of Two Commutative Binary Operations on
$\Z/10\Z$}, manuscript in preparation, 2026.

\bibitem{KatokUgarcovici2007} S.~Katok, I.~Ugarcovici,
\emph{Symbolic dynamics for the modular surface},
Bull.\ AMS \textbf{44} (2007), 87--132.

\bibitem{Morishita2024} M.~Morishita,
\emph{Knots and Primes}, 2nd ed., Springer, 2024.

\bibitem{MatsusakaUeki2023} T.~Matsusaka, J.~Ueki,
\emph{Rademacher symbols for triangle groups}, 2023.

\bibitem{Burrin2024} C.~Burrin, F.~von Essen,
\emph{Cusp winding and the Rademacher symbol}, 2024.

\bibitem{Smith2007} J.~D.~H.~Smith,
\emph{An Introduction to Quasigroups and Their Representations},
Chapman \& Hall / CRC, 2007.

\bibitem{Pflugfelder1990} H.~O.~Pflugfelder,
\emph{Quasigroups and Loops: Introduction}, Heldermann, 1990.

\end{thebibliography}

\end{document}
